\documentclass[14pt, reqno]{extarticle}

\usepackage{amsmath, amsthm, amssymb}
\usepackage{enumitem}
\usepackage{tcolorbox}
\usepackage{hyperref}
\usepackage{tikz}
\usepackage{tikz-cd}
\usepackage{pgfplots}
\pgfplotsset{compat=1.18}
\usetikzlibrary{arrows.meta}
\usepackage{mathrsfs}
\usepackage{fancyhdr}
\usepackage[bottom=0.75in, top=1in, left=0.5in, right=0.5in]{geometry}
\usepackage{array}   % for \newcolumntype macro
\newcolumntype{L}{>{$}l<{$}}



\theoremstyle{plain}
\newtheorem{theorem}{Theorem}[section]
\newtheorem{proposition}{Proposition}[section]
\newtheorem{exercise}{Exercise}
\newtheorem{lemma}{Lemma}[section]
\newtheorem{corollary}{Corollary}[section]
\theoremstyle{definition}
\newtheorem{definition}{Definition}[section]
\newtheorem*{example}{Example}

% Pretty good font with math typesetting
% ~~~~~~~~~~~~~~~~~~~~~~~~~~~~~~~~~~~~~~~~~~~
\usepackage{concmath}
\usepackage[T1]{fontenc}

% funny and thin
% \renewcommand*\rmdefault{cmdh}
% \usepackage[T1]{fontenc}

%very good and sciency
% \usepackage[sfdefault]{ClearSans} %% option 'sfdefault' activates Clear Sans as the default text font
% \usepackage[T1]{fontenc}
% \usepackage{euler}

\theoremstyle{remark}
\newtheorem*{remark}{Remark}

\renewcommand{\phi}{\varphi}
\renewcommand{\epsilon}{\varepsilon}
\renewcommand{\emptyset}{\varnothing}

\newcommand{\RR}{\mathbb{R}}
\newcommand{\ZZ}{\mathbb{Z}}
\newcommand{\NN}{\mathbb{N}}
\newcommand{\CC}{\mathbb{C}}
\newcommand{\QQ}{\mathbb{Q}}

\newcommand{\der}[2]{\frac{\partial #1}{\partial #2}}

\DeclareMathOperator{\ima}{\text{im}}

\begin{document}

\title{Pseudodifferential Operators}
\author{Henry Woodburn}
\maketitle

\section{Differential Operators}

We would like to generalize the notion of a differential
operator on $C^\infty(R^n)$ to sections of vector bundles on a manifold. Consider the 
following system of equations:
\[
\begin{matrix}
    3 \frac{\partial}{\partial x} f(x,y,z) + \frac{\partial}{\partial x} g(x,y,z) 
    + 2 \frac{\partial}{\partial z} h(x,y,z) & = 0\\
    3 \frac{\partial}{\partial x} f(x,y,z) + 3 \frac{\partial}{\partial x} f(x,y,z)
    + 3 f(x,y,z) & = 0
\end{matrix}
\]

We can represent the right side as an operator $D$ applied to an $R^3$ valued function,
where
\[
    D = \begin{bmatrix} 3 & 1 & 0\\ 3 & 1 & 0\end{bmatrix}\frac{\partial}{\partial x}
    + \begin{bmatrix} 0 & 0 & 1 \\ 0 & 0 & 0\end{bmatrix}\frac{\partial}{\partial z}
    + \begin{bmatrix} 0 & 0 & 0 \\ 0 & 0 & 3\end{bmatrix}.
\]

In other words, $D$ is a linear map $C^\infty(\RR^2, \RR^3) \to C^\infty(\RR^2, \RR^2)$. 
This suggests that instead of considering differential operators acting on functions 
(sections of a trivial line bundle,) we consider differential operators which 
act on sections of vector bundles. 

In the following section, we let $X$ be a manifold, with $E$ and $F$ smooth complex
vector bundles on $X$. We denote by $\Gamma(E)$ the space of smooth $C^\infty$ sections of $E$. 

\begin{definition}\label{diff}
    A \textit{differential operator} of order $m$ on $X$ is a 
    linear map $P: \Gamma(E) \to \Gamma(F)$, such that if we choose local
    coordinates $x_i$ in an open set $U \subset X$, and also local trivializations
    $E|_U \simeq U \times \CC^p$ and $F|_U \simeq U \times \CC^q$, then $P$ has the form
    \[
        P = \sum_{|\alpha| \leq m} A^\alpha(x) \frac{\partial^{|\alpha|}}{\partial x^\alpha}
    \]
    where the $A^\alpha$'s are smooth $p\times q$ matrices of smooth functions, and
    such that one of the highest order coefficients is nonzero somewhere. 
\end{definition}

Observe that if we make a change of local trivializations of $E|_U$ and $F|_U$ by maps
$g_E: U \to GL_p(\CC)$ and $g_F: U \to GL_p(\CC)$, then in these trivializations,
$P$ has the form
\[
    P = g_F\left(\sum_{|\alpha| \leq m} A^{\alpha} \frac{\partial^{|\alpha}}{\partial x^\alpha}\right)g_E^{-1}
    = \sum_{|\alpha| \leq m} \tilde{A}^\alpha \frac{\partial^{|\alpha|}}{\partial x^\alpha},
\]

where the coefficients $\tilde{A}^\alpha$ can be obtained using the product rule 
to interchange the orders of $\partial_x$ and $g_E^{-1}$. Namely, for all $|\alpha| = m$, 
we have the formula 
\[
    \tilde{A}^\alpha = g_E A^\alpha g_E^{-1}
\]
for the modified coefficients. 

Similarly, if we change to coordinates $\tilde{x} = \tilde{x}(x)$ on the set $U$, $P$ takes the form
\[
    P = \sum_{|\alpha| \leq m} \tilde{A}^\alpha(\tilde{x})\frac{\partial^{|\alpha|}}{\partial x^\alpha}.
\]

Using the Jacobian for the change of variables, for any $j$ we can write
\[
    \frac{\partial}{\partial x_j} = \sum_{k = 1}^n \frac{\partial \tilde{x}_k}{\partial x_j} \frac{\partial}{\partial \tilde{x}_k}.
\]

For any multi-index $\alpha$ with $|\alpha| = m$, we can write 
\[
    \frac{\partial}{\partial x^\alpha} = \sum_{|\beta| = m} \left[\frac{\partial\tilde{x}}{\partial x}\right]_\beta^\alpha \frac{\partial}{\partial \tilde{x}^\beta},
\]
where $\left[\frac{\partial\tilde{x}}{\partial x}\right]_*^*$ denotes the symmetrization of the $m$-th exterior 
power of the jacobian. Note that both upper and lower indices are symmetrized to reflect the commutativity 
of partial derivatives. 

Thus we have 
\[
    \tilde{A}^\alpha = \sum_{|\beta| = m} A^\beta \left[ \frac{\partial{\tilde{x}}}{\partial x}\right]^\alpha_\beta
\]
for the top order coefficients with $|\alpha| = m$. 

We see that $P$ behaves well under coordinate changes and changes of trivialization, so long as we only 
consider the top order terms. In fact, we have proven that the coefficients $\{i^m A^\alpha\}_{|\alpha| = m}$
represent a well defined section $\sigma(P)$ of the bundle $(\odot^m TX) \otimes \text{Hom}(E,F)$, where
$\odot$ denotes the symmetric tensor product. 

Next we note that the space $\odot^m V$ is \textit{canonically} isomorphic to the space of homogeneous polynomial
functions of degree $m$ on $V^*$. Hence, for each cotangent vector $\xi \in T_x^* X$, the principal 
symbol gives a linear map 
\[
    \sigma_\xi(P): E_x \to F_x.
\]
In other words, the symbol can be considered as a (polynomial) function on the cotangent bundle. 



\subsection{Examples on Riemannian Manifolds}

Now let $X$ be a riemannian manifold equipped with a metric $\langle\cdot,\cdot\rangle$. In local 
coordinates we can express the metric using the basis as
\[
    \sum g_{jk}dx_j dx_k
\]
for some matrix $g_{jk}$. We also denote by $g^{jk}$ the inverse of $g_{jk}$, and $g := \det(\{g_{jk}\})$.

\begin{example}
    Consider the Laplace-Beltrami operator $\Delta: C^\infty(X) \to C^\infty(X)$ defined by $\Delta = d^*d$, where
    $d^*$ is defined on $k$-forms $\alpha$ by $d^*\alpha = (-1)^{(nk+n+1)}*d^*\alpha$. We first calculate 
    an explicit formula for $d^*$ applied to $k$-forms $\alpha = \sum A_i dx^i$:
    \[
        d^*\alpha = \frac{1}{\sqrt{g}}\sum_{ij}\frac{\partial}{\partial x_j}(A_i g^{ij}\sqrt{g}).
    \]
    Then we can apply it to $\alpha = df = \sum \frac{\partial f}{\partial x_i} dx^i$ to obtain the following formula for $\Delta$:
    \begin{align}
        \Delta f & = \frac{1}{\sqrt{g}}\sum_{ij}\frac{\partial}{\partial x_j}(\sqrt{g} g^{ij} \frac{\partial f}{\partial x_i} )\\
        & = \sum_{j,k=1}^n g^{jk} \frac{\partial^2 f}{\partial x_j \partial x_k} + \{\text{lower order terms}\}.
    \end{align}

    Now we can calculate its symbol: for some cotangent vector $\xi = \sum \xi_k dx^k$, the polynomial
    function is given by
    \[
        \sigma_\xi(\Delta) = i^2 \sum_{j,k = 1}^n g^{jk} \xi_j \xi_k = -\|\xi\|^2.
    \]
\end{example}

\begin{example}
    From this discussion, it is natural to wonder if we can just define an operator with symbol $i\|\xi\|$. 
    In fact we could, but it will be pseudodifferential and not quite what we are after in several ways. 

    Instead, we introduce the concept of clifford action of $T^*M$ on a bundle $S$. In defining the 
    dirac operator $D$ in the usual way, we see that its symbol is 
    \[
        \sigma_\xi(D) = ic(\xi),
    \]
    where $c(\xi) \in \text{End},mk(S)$ denotes clifford multiplication by $\xi$. Trivializing $S$ in some neighborhood $U$, $c(\xi)$ becomes 
    some matrix-valued function on $U$, and thus $D$ is an ordinary differential operator on $S$. 
\end{example}

\begin{remark}
    For a differential operator $P: \Sigma(E) \to \Sigma(F)$ we can consider $\sigma(P)$ as a bundle map of the pullbacks 
    $\pi^*E \to \pi^*F$ in the natural way, under the projection $\pi: T^*M \to M$.
\end{remark}

\section{Sobolev Spaces}

The point of this section is to introduce the fourier transform-based sobolev spaces $H^s$ for $s \in \RR$,
and to extend them to sections of vector bundles in a suitable way. 

\begin{definition}\label{orig}
    Let $E$ be a hermitian vector bundle with connection $\nabla$ on a compact Riemannian manifold $X$. Given 
    $u \in \Sigma(E)$, we have $\nabla u \in \sigma(T^*X \otimes E)$. Furthermore, using the tensor product
    connection on $T^*X \otimes E$, $\nabla \nabla u \in \sigma(T^*X\otimes T^*X \otimes E)$. Thus
    for any $k$, we can define a norm
    \[
        \|u\|_k^2 = \sum_{j = 0}^k \int_X |\nabla^k u|,
    \]
    called the basic Sobolev $k$-norm on $\sigma(E)$.

    The completion of $\sigma(E)$ under this norm is the Sobolev space $H^k(E)$.
\end{definition}

A fundamental fact about sobolev spaces is the following proposition:

\begin{proposition}\label{basic}
    A differential operator $P: \sigma(E) \to \sigma(E)$ of order $m$ is bounded 
    as a linear map $P: H^k \to H^{k-m}$. Hence, $P$ extends to a linear
    map $P: H^k \to H^{k-m}$. 
\end{proposition}

The next construction reduces the study of sections of $E$ to the study of smooth $\CC^p$-valued functions
with compact support in the unit ball using a partition of unity. Then we can build a norm on $\sigma(E)$ using
the fourier transform on sections in each trivialized patch of $E$ under the transformed coordinates.
I will not provide it here, but we need an open cover of X consisting of trivialized patches of $E$, 
and a partition of unity subordinate to this cover, such that the open cover is sufficiently overlapping.
Then we change coordinates with a bijection $[-1,1] \to \RR$. 

\subsection{Fourier transform-based Sobolev spaces}

Here we briefly outline the sobolev spaces $H^s(\RR^n)$, with values in $\CC^p$, for $s \in \RR$. 
These are constructed using the fourier transform.

\begin{definition}
    The \textit{Fourier Transform} $\hat{u}: \RR^n \to \CC^n$ 
    of a compactly supported function $u \in C^\infty(\RR^n, \CC^n)$ 
    is defined as 
    \[
        \hat{u}(\xi) = (2\pi)^{-n/2} \int_{\RR^n} e^{-i\langle x, \xi\rangle}u(x)dx.
    \]
    Denote the Fourier map by $\mathcal{F}$.
\end{definition}

In fact, $\mathcal{F}$ is an isometry $\mathscr{S} \to \mathscr{S}$, where $\mathscr{S}$ 
is the space of \textit{Schwartz functions} of ``rapidly decreasing functions,'' 
who are dominated by polynomials of arbitrarily large negative order on any compact set,
under the $L^2$ norm. 

\begin{definition}
    We define the \textit{Sobolev $s$-norm} $\|\cdot\|_s$ by the formula
    \[
        \|u\|_s^2 = \int (1 + \|\xi\|)^{2s}|\hat{u}(\xi)|^2 d\xi.
    \]
\end{definition}

To motivate this definition, note that for compactly supported $u$ on $\RR$, after integrating by parts, we have
\[
    \widehat{(\partial_{x}u)} = \int e^{ix\xi} \partial_{x}(u)(x)dx
    = \int \xi e^{ix\xi} u(x) dx = \xi \hat{u}.
\]
Thus for the case $s = 1$, we have
\[
    \int |\xi|^2|\hat{u}(\xi)|^2d\xi = \int {|\xi\hat{u}(\xi)|}^2d\xi = \int {|\widehat{(\partial_x u)}(\xi)|}^2d\xi 
    = \int {|\partial_x u(x)|}^2 dx
\]
since $\mathcal{F}$ is an isometry in the $L^2$ norm. 

One surprising result is the following:

\begin{proposition}[The Sobolev Embedding Theorem]
    Fix a positive integer $k$. Then for each real number $s > (n/2) + k$, there is a constant
    $K_s$ such that for every $u \in \mathscr{S}$, 
    \[
        \|u\|_{C^k} \leq K_s \|u\|_s.
    \]
    Hence, there is a continuous embedding 
    \[
        H^s \subset C^k.
    \]
\end{proposition}

In order to extend these results globally to define Sobolev spaces on compact manifolds, 
we need a couple of lemmas, which we will not prove, but are not too unreasonable. 

\begin{lemma}\label{mtx}
    Let $A$ be a smooth matrix-valued function on $\RR^n$ so that $|D^\alpha A|$ is bounded for all $\alpha$. 
    Then the map $T: \mathscr{S} \to \mathscr{S}$ given by $Tu = Au$ extends to a bounded
    linear map $T: H^s \to H^s$ for all $s \in \RR$. 
\end{lemma}

\begin{lemma}\label{btx}
    Let $\Phi: \RR^n \to \RR^n$ be a $C^\infty$ diffeomorphism which is linear outside
    a compact set. Then the map $T: \mathscr{S} \to \mathscr{S}$ given by $Tu = u \circ \Phi$
    extends to a bounded linear map $T: H^s \to H^s$ for all $s \in \RR$. 
\end{lemma}

Now we globalize. Let $E$ be a smooth vector bundle over a compact manifold $X$. Fix a good presentation
$\{U_\beta\}$ for $E$ with coordinates $x_\beta: U_\beta \to \RR^n$ and a partition of 
unity $\{\chi_\beta\}$ subordinate to the cover $\{B_\beta\} = \{x^{-1}_\beta(B^n)\}$. 

Since $\sum \chi_\beta = 1$, any $u \in \Gamma(E)$ can be written $u = \sum u_\beta$ where
$u_\beta = \chi_\beta u$. For $s \in \RR$, define a Sobolev $s$-norm on $\Gamma(E)$ by
setting
\[
    \|u\|_s = \sum_{\beta = 1}^N \|u_\beta\|_s,
\]
where the $\|u_\beta\|_s$ are defined in terms of the presentation $E|_{U_\beta} \simeq \RR^n \times \CC^n$.

As a consequence of \ref{mtx} and \ref{btx}, we have the following: 

\begin{proposition}
    For any $s \in \RR$ and any smooth vector bundle $E$ over a compact manifold, the equivalence class
    of the norm $\|\cdot\|_s$ as defined above is independent of the good presentation used to define it.
    Furthermore, when $s = k$ is an integer, the equivalence class of $\|\cdot\|_k$ is the same
    as that defined in the first definition \ref{orig}.
\end{proposition}

Moreover, we have an analagous result to \ref{basic} for differential operators:

\begin{proposition}
    Let $P: \Gamma(E) \to \Gamma(F)$ be a differential operator of order $m$ as defined 
    in \label{diff}
    whose (locally expressed) coefficients are bounded as in \ref{mtx}. Then $P: \Gamma(E) \to \Gamma(E)$ 
    extends to a bounded linear map $P: H^s(E) \to H^{s-m}(F)$ for all $s \in \RR$. 
\end{proposition}

\begin{remark}
    In fact, $H^s(E)$ is a hilbert space under the $H^s$ inner product
    \[
        \langle u, v\rangle_s = \int {(1 + |\xi|)}^{2s}\hat{u}(\xi)\overline{\hat{v}(\xi)}d\xi.
    \]

    Alternatively, we can define a pairing 
    $\langle\cdot,\cdot\rangle_*: H^s(E) \times H^{-s}(E^*) \to \RR$ by setting 
    \[
        \langle u, v\rangle_* = \int u(v),
    \]
    placing $H^s(E)$ and $H^{-s}(E^*)$ in duality. 
\end{remark}

\end{document}